\documentclass[11pt]{article}

\usepackage[T1]{fontenc}
\usepackage[utf8]{inputenc}
\usepackage{lmodern}
\usepackage[margin=1in]{geometry}
\usepackage{microtype}
\usepackage{amsmath,amssymb,mathtools}
\usepackage{booktabs,tabularx,array}
\usepackage{enumitem}
\usepackage{xcolor}
\usepackage{hyperref}

\hypersetup{
  colorlinks=true,
  linkcolor=blue!55!black,
  citecolor=blue!55!black,
  urlcolor=blue!55!black
}

\newcommand{\Aphys}{\mathcal{A}_{\mathrm{phys}}}
\newcommand{\R}{\mathbb{R}}
\newcommand{\cO}{\mathcal{O}}
\newcommand{\Spec}{\operatorname{Spec}}
\newcommand{\Tr}{\operatorname{Tr}}
\newcommand{\YM}{\mathrm{YM}}

\title{Toward a Phenomenon-Preserving Calculus\\
of Physical Representations\\
\large The Yang--Mills Mass Gap as a Problem of Representational Transport}
\author{}
\date{}

\begin{document}

\maketitle

\begin{abstract}
The fragmented perspectives on the Yang--Mills mass-gap problem point toward a
need deeper than another choice of field variables.  What is missing is a
mathematics of the relations between representations: a framework that states
which physical phenomena, structural conditions, and quantitative bounds are
preserved when variables, regulators, scales, or apparent ontologies change.
This essay develops a phenomenon-first conception of physical representation.
It separates defining physical content, admissibility conditions, emergent
targets, and auxiliary constructs; introduces representation packages and
proof-bearing morphisms; and treats the mass gap as infrared spectral rigidity
that must be generated in one representation and transported to the physical
continuum spectrum.  The resulting perspective reframes the search for a mass
gap as the search for a representation in which non-Abelian dynamics becomes a
uniform positive inequality together with certified transformations carrying
that inequality to the reconstructed quantum theory.
\end{abstract}

\section{The new representational need}

The fragmented descriptions of Yang--Mills theory suggest a central conclusion:

\begin{quote}
We do not merely need a better representation of Yang--Mills fields.  We need a
mathematical representation of the relations between representations---one
that records which physical phenomena survive when variables, scales,
regulators, and even apparent ontologies change.
\end{quote}

The classical action
\begin{equation}
  S[A]=\frac{1}{4g^2}\int_{\R^4}
  \Tr(F_{\mu\nu}F^{\mu\nu})\,d^4x,
  \qquad F=dA+A\wedge A,
  \label{eq:action}
\end{equation}
is an exceptionally economical representation of local geometric dynamics.
It is not, however, a representation in which the expected infrared structure
is manifest.  The gauge potential is redundant, the classical action is
scale-free, and perturbative excitations are massless, whereas the expected
physical theory has a dynamically generated scale and a gauge-invariant
massive spectrum.

Representing a theory by its preferred variables and equations therefore risks
identifying a descriptive instrument with physical ontology.  The mass-gap
problem instead encourages the shift
\begin{equation}
  \text{representation of objects}
  \quad\longrightarrow\quad
  \text{representation of physical invariants and their transformations}.
  \label{eq:shift}
\end{equation}

\section{Why the usual notion of representation is insufficient}

A conventional formulation normally selects fundamental variables, equations
of motion, a state space, observables, and a quantization procedure.  Beginning
with the connection $A_\mu$ already carries several tacit commitments:
\begin{itemize}[leftmargin=2em]
  \item the gauge potential is treated as the natural primitive object;
  \item locality is expected to be manifest at the level of primitive
        variables;
  \item the classical configuration space is presumed to remain the right
        ontology after quantization;
  \item the quantum theory is approached as a deformation of the classical
        theory;
  \item the particles visible around a perturbative vacuum are treated as
        suitable coordinates for the infrared regime.
\end{itemize}

These commitments need not all be false.  The conceptual problem is that they
may be poorly aligned with the target phenomenon.  A representation can be
faithful and calculationally powerful in one regime while being structurally
blind to another.  In Yang--Mills theory the classical and perturbative
representations make the input law transparent while rendering the emergent
infrared conclusion obscure.

\section{A phenomenon ledger}

The analysis should begin by separating four epistemic classes: defining
physical content, structural admissibility conditions, target phenomena, and
associated but logically distinct phenomena.

\subsection{Defining physical content}

The following features identify the intended theory or its physical
universality class:
\begin{itemize}[leftmargin=2em]
  \item four-dimensional continuum spacetime in the final theory;
  \item a compact simple gauge group $G$;
  \item the non-Abelian short-distance interaction structure;
  \item asymptotic freedom and its ultraviolet scaling;
  \item a sufficiently rich algebra of gauge-invariant observables and their
        correlation functions;
  \item nontrivial interaction.
\end{itemize}
Discarding these would risk solving a different problem.

\subsection{Structural admissibility conditions}

These conditions determine whether the constructed object is an acceptable
quantum field theory:
\begin{itemize}[leftmargin=2em]
  \item positivity of physical probabilities;
  \item unitarity after reconstruction;
  \item Poincar\'e covariance in the continuum;
  \item causal locality for physical observables;
  \item a stable vacuum and positive energy;
  \item compatibility of Euclidean and Hamiltonian descriptions;
  \item controlled continuum and infinite-volume limits.
\end{itemize}
They need not be manifest in every intermediate representation, but every
temporary loss requires a recovery theorem.

\subsection{Target phenomena}

The following are conclusions to be generated rather than assumptions to be
preserved:
\begin{itemize}[leftmargin=2em]
  \item dimensional transmutation;
  \item a finite physical correlation length;
  \item exponential clustering;
  \item a positive spectral gap;
  \item a spectrum of massive gauge-invariant excitations.
\end{itemize}
This distinction guards against circularity.  A construction that assumes a
uniform exponential-decay estimate has already incorporated much of the desired
mass-gap conclusion.

\subsection{Associated but separate phenomena}

Confinement, a Wilson-loop area law, flux-tube formation, detailed glueball
masses, and topological susceptibility may help identify the dynamical
mechanism.  They are not interchangeable with the exact mass-gap statement.
The logical implications among them require separate theorems.

\section{Which constructs may be compromised?}

A phenomenon-first approach distinguishes physical indispensability from
historical familiarity.

\begin{table}[ht]
\centering
\small
\renewcommand{\arraystretch}{1.22}
\begin{tabularx}{\textwidth}{>{\raggedright\arraybackslash}p{0.25\textwidth}
>{\raggedright\arraybackslash}p{0.20\textwidth}X}
\toprule
\textbf{Construct} & \textbf{Representational status} &
\textbf{Physical content to preserve} \\
\midrule
Gauge potential $A_\mu$
  & Potentially replaceable
  & Gauge-invariant observable relations \\
Particular gauge fixing
  & Auxiliary
  & Independence of physical predictions from the gauge \\
Ghost fields
  & Representation-dependent
  & Correct physical state space and unitarity \\
Perturbative Fock vacuum
  & Potentially inadequate in the infrared
  & Existence and stability of the physical vacuum \\
Gluons as asymptotic particles
  & Not required in pure Yang--Mills theory
  & Gauge-invariant physical spectrum \\
Classical configurations as fundamental ontology
  & Potentially replaceable
  & Quantum observables and correlation structure \\
Manifest locality of primitive variables
  & Potentially relaxable
  & Causal locality of physical observables \\
Smooth continuum at intermediate stages
  & Regulatable
  & Exact recovery of the continuum limit \\
Lorentz symmetry at finite cutoff
  & Temporarily relaxable
  & Poincar\'e covariance in the limit \\
Path integral over $A$
  & One construction strategy
  & A positive consistent state on physical observables \\
Wilson loops as privileged primitives
  & Useful but not sacrosanct
  & The gauge-invariant information they encode \\
Specific renormalization scheme
  & Auxiliary
  & Universal scaling and scheme-independent physics \\
Particle ontology
  & Emergent
  & Spectral and scattering content where defined \\
\bottomrule
\end{tabularx}
\caption{A construct-compromise analysis for Yang--Mills representations.}
\label{tab:compromise}
\end{table}

The central distinction is between \emph{manifest} and \emph{recovered}
structure.  Locality may fail to be manifest in a dual variable while causal
locality remains exactly reconstructible for physical observables.  A lattice
regulator may break continuous spacetime symmetry while the limiting theory
recovers Poincar\'e covariance.  The gauge potential can lose its status as a
primitive without abandoning the gauge-invariant content of Yang--Mills theory.

\section{Physical reality as an invariant intersection}

The physical theory should not be identified with any one familiar
representation.  It may instead be conceived as the structure invariant across
an atlas:
\begin{equation}
\begin{aligned}
  \mathfrak R_{\mathrm{lattice}}
  &\longleftrightarrow
  \mathfrak R_{\mathrm{Euclidean}}
  \longleftrightarrow
  \mathfrak R_{\mathrm{observable}}\\
  &\longleftrightarrow
  \mathfrak R_{\mathrm{Hamiltonian}}
  \longleftrightarrow
  \mathfrak R_{\mathrm{RG}}
  \longleftrightarrow
  \mathfrak R_{\mathrm{dual}}.
\end{aligned}
\label{eq:atlas}
\end{equation}

The analogy with coordinate charts on a manifold is helpful but incomplete.
Transformations between physical representations can be singular, exact only
in a limit, valid only within a scale regime, nonlocal, or irreversible under
coarse-graining.  They may change the apparent ontology and may preserve some
observables while discarding others.  Representations alone are therefore not
enough: the transformations between them must become proof-bearing mathematical
objects.

\section{A phenomenon-preserving representation package}

A representation may be organized as a typed package
\begin{equation}
  \mathfrak R
  =
  \bigl(
    \mathcal K,
    \Aphys,
    \mathcal S,
    \mathcal D,
    \mathcal G,
    \mathcal L,
    \mathcal Q
  \bigr),
  \label{eq:package}
\end{equation}
where:
\begin{description}[leftmargin=2.8cm,style=nextline]
  \item[$\mathcal K$] is the carrier or kinematic presentation: connections,
    lattice links, loops, operators, or dual variables;
  \item[$\Aphys$] is the algebra or designated class of physical observables;
  \item[$\mathcal S$] is the admissible state space;
  \item[$\mathcal D$] specifies dynamics, correlations, or time evolution;
  \item[$\mathcal G$] records physical symmetries, redundancies, equivalence
    relations, and constraints;
  \item[$\mathcal L$] specifies scale transformations and limiting operations;
  \item[$\mathcal Q$] contains quantitative control data: norms, errors, decay
    rates, and domains of uniformity.
\end{description}

The separation of $\mathcal K$ from $\Aphys$ is decisive.  Two representations
can have radically different carriers while encoding the same observable
theory.  Conversely, a formal change of variables that resembles an algebraic
equivalence need not be physically adequate if it fails to preserve states,
positivity, or the limiting structure.

\section{Proof-bearing morphisms}

Let
\begin{equation}
  \Phi:\mathfrak R\longrightarrow\mathfrak R'
\end{equation}
be a transformation between representations.  It should include an observable
pullback
\begin{equation}
  \Phi^*:\Aphys'\longrightarrow\Aphys
\end{equation}
and an appropriate transformation of states.  Exact physical equivalence would
require
\begin{equation}
  \omega'(O')=\omega\!\left(\Phi^*O'\right)
  \label{eq:exact-equivalence}
\end{equation}
for the relevant states and observables.

For a regulated, approximate, or scale-dependent translation, the morphism
must instead supply a quantitative certificate
\begin{equation}
  \left|
    \omega'(O')-\omega\!\left(\Phi^*O'\right)
  \right|
  \leq
  \varepsilon_\Phi(O';a,L,\mu).
  \label{eq:approx-equivalence}
\end{equation}
The decisive question is whether
\begin{equation}
  \sup_{O'\in\mathcal C}
  \varepsilon_\Phi(O';a,L,\mu)
  \longrightarrow0
  \label{eq:uniform-error}
\end{equation}
for a convergence-determining class $\mathcal C$, with the uniformity required
for the continuum and infinite-volume limits.

Every morphism should state its epistemic type: exact equivalence, gauge
quotient, regularization, coarse-graining, reconstruction, continuum limit,
asymptotic equivalence, duality, or heuristic correspondence.  This prevents a
numerical approximation, an effective theory, a conjectural duality, and an
exact equivalence from being treated as though they carried the same force.

\section{Phenomenon-preservation certificates}

Transformations should carry explicit proof obligations.

\paragraph{Observable preservation.}
Does the morphism preserve expectation values and correlation functions for a
sufficiently rich gauge-invariant observable class?

\paragraph{Positivity preservation.}
Does a positive Euclidean functional remain capable of reconstructing a
positive physical Hilbert space?

\paragraph{Symmetry preservation.}
Which symmetries are exact, broken by the regulator, anomalous, or recovered
only in the limit?

\paragraph{Locality preservation.}
Is locality manifest, weakened, or reconstructed?  Can nonlocal primitive
variables still yield causal locality of physical observables?

\paragraph{Scale preservation.}
Does the transformation retain the correct beta function, generated scale,
and universality class?

\paragraph{Limit preservation.}
Are the bounds uniform in lattice spacing, volume, and observation distance?
Can the relevant limiting operations be exchanged?

\paragraph{Spectral preservation.}
Does a decay, mixing, coercivity, or contraction estimate in one representation
imply a positive spectral bound after physical reconstruction?

The last obligation is the core of the mass-gap problem.

\section{The gap across representations}

The same expected phenomenon has distinct mathematical manifestations.

\begin{table}[ht]
\centering
\small
\renewcommand{\arraystretch}{1.2}
\begin{tabularx}{\textwidth}{>{\raggedright\arraybackslash}p{0.29\textwidth}X}
\toprule
\textbf{Representation} & \textbf{Candidate gap witness} \\
\midrule
Hamiltonian
  & $H\geq\Delta(1-P_0)$ \\
Euclidean correlations
  & Uniform exponential clustering \\
Transfer operator
  & Separation between leading and subleading spectral values \\
Probability measure
  & Uniform mixing, Poincar\'e inequality, or log-Sobolev inequality \\
Renormalization group
  & Controlled flow into a massive regime with finite physical correlation length \\
Effective action
  & Quantitative coercivity around the physical vacuum \\
Observable algebra
  & Absence of energy-momentum spectral support near zero \\
Lattice theory
  & A gap estimate uniform in volume and positive in physical units as $a\to0$ \\
\bottomrule
\end{tabularx}
\caption{Representation-dependent witnesses of infrared rigidity.}
\label{tab:gap-witnesses}
\end{table}

These witnesses are not automatically equivalent.  The new framework must
contain transfer theorems of the form
\begin{equation}
  \text{certificate in }\mathfrak R_i
  \quad\Longrightarrow\quad
  \text{mass gap in }\mathfrak R_{\mathrm{physical}}.
  \label{eq:transfer-pattern}
\end{equation}
One quantitative form would be
\begin{equation}
  \Delta_{\mathrm{physical}}
  \geq c_i g_i-\varepsilon_i,
  \qquad c_i>0,
  \label{eq:gap-bound}
\end{equation}
where $g_i$ is a representation-dependent rigidity witness and
$\varepsilon_i$ vanishes uniformly through the required limits.

\section{From emergence to coercivity}

In Hamiltonian language, the mass gap is a universal exclusion statement: no
unknown sequence of physical states may have positive energies tending to
zero.  A more favorable representation could turn this negative statement into
a positive inequality.  Candidate forms include:
\begin{itemize}[leftmargin=2em]
  \item a Poincar\'e or logarithmic Sobolev inequality;
  \item contraction of a transfer or coarse-graining operator;
  \item uniform probabilistic mixing;
  \item coercivity of a gauge-invariant effective action;
  \item a monotone renormalization-group quantity;
  \item a direct operator inequality above the vacuum sector.
\end{itemize}

For example, one might seek
\begin{equation}
  \operatorname{Var}_{\mu}(f)
  \leq
  \frac{1}{\lambda}\mathcal E_{\mu}(f,f),
  \qquad \lambda>0,
  \label{eq:poincare}
\end{equation}
and a theorem transporting $\lambda$ to the physical spectral gap.  Such a
reformulation would turn collective emergence into mathematical rigidity.  Its
value depends entirely on proving that the inequality is gauge invariant,
nonperturbatively valid, uniform in volume, stable under renormalization,
correctly scaled in the continuum limit, and sufficient for spectral
reconstruction.

\section{A hierarchy of compromise and preservation}

The framework imposes the following discipline:
\begin{equation}
\boxed{
\begin{array}{c}
\textbf{May change freely}\\[1mm]
\text{variables, gauges, regulators, coordinates, auxiliary fields}\\[3mm]
\downarrow\\[-1mm]
\textbf{May change with a recovery theorem}\\[1mm]
\text{manifest locality, Lorentz symmetry, continuum form, particle language}\\[3mm]
\downarrow\\[-1mm]
\textbf{Must be preserved}\\[1mm]
\text{physical observables, positivity, causality, universality class}\\[3mm]
\downarrow\\[-1mm]
\textbf{Must be derived}\\[1mm]
\text{dynamical scale, exponential clustering, positive mass gap}
\end{array}}
\label{eq:hierarchy}
\end{equation}

This hierarchy guards against opposite errors.  Preserving every familiar
construct may trap the proof in a representation that hides the phenomenon.
Relaxing defining physical structure may produce a gap for a different theory.

\section{The deeper conceptual lesson}

The physical content of Yang--Mills theory may reside neither in fields,
particles, loops, nor lattice variables considered in isolation.  It may reside
in a network of stable relations across representations.  On this view,
\begin{itemize}[leftmargin=2em]
  \item gauge fields are coordinates on physical information;
  \item regulators are temporary realizations;
  \item particles are spectral interpretations;
  \item renormalization is a transformation of descriptive scale;
  \item the continuum theory is a coherence condition across limits;
  \item the mass gap is an invariant form of infrared rigidity.
\end{itemize}

The missing theory is thus a form of \emph{representational dynamics}: a
mathematics of how physical descriptions transform, what they forget, what
they preserve, and which emergent properties become coercive in each one.

\section{Central proposition}

The research programme can be organized around the following proposition:

\begin{quote}
The Yang--Mills mass-gap problem is fundamentally the problem of generating and
transporting infrared spectral rigidity across a chain of representations whose
natural variables separately expose locality, gauge invariance, positivity,
scale dependence, and spectral structure.
\end{quote}

This changes the search question from
\begin{quote}
Which existing formalism can directly prove the mass gap?
\end{quote}
to
\begin{quote}
In which representation does non-Abelian quantum dynamics become a uniform
positive inequality, and through which certified transformations can that
inequality be transported to the physical continuum spectrum?
\end{quote}

The missing object may therefore be neither a new field variable nor a new
computational technique alone.  It may be a phenomenon-preserving calculus of
representations capable of turning collective quantum emergence into
transferable mathematical rigidity.

\begin{thebibliography}{99}

\bibitem{jaffe-witten}
A.~Jaffe and E.~Witten,
\emph{Quantum Yang--Mills Theory},
in \emph{The Millennium Prize Problems},
Clay Mathematics Institute and American Mathematical Society, 2006.

\bibitem{osterwalder-schrader}
K.~Osterwalder and R.~Schrader,
\emph{Axioms for Euclidean Green's Functions},
Communications in Mathematical Physics \textbf{31} (1973), 83--112;
\textbf{42} (1975), 281--305.

\bibitem{wilson}
K.~G.~Wilson,
\emph{Confinement of Quarks},
Physical Review D \textbf{10} (1974), 2445--2459.

\bibitem{gribov}
V.~N.~Gribov,
\emph{Quantization of Non-Abelian Gauge Theories},
Nuclear Physics B \textbf{139} (1978), 1--19.

\bibitem{haag}
R.~Haag,
\emph{Local Quantum Physics: Fields, Particles, Algebras},
Springer, second edition, 1996.

\end{thebibliography}

\end{document}
